\documentclass[10pt]{article}
\usepackage[a4paper,margin=0.72in]{geometry}
\usepackage{array}
\usepackage{amsmath}
\usepackage{booktabs}
\usepackage{enumitem}
\usepackage{fontspec}
\usepackage{hyperref}
\usepackage[table]{xcolor}
\usepackage{titlesec}

\hypersetup{colorlinks=true,linkcolor=black,urlcolor=blue,citecolor=black}
\setlength{\parindent}{0pt}
\setlength{\parskip}{3pt}
\setlist[itemize]{leftmargin=1.1em,itemsep=1pt,topsep=2pt}
\setlist[enumerate]{leftmargin=1.25em,itemsep=1pt,topsep=2pt}
\renewcommand{\arraystretch}{1.13}
\titleformat{\section}{\large\bfseries}{\thesection}{0.45em}{}
\titlespacing*{\section}{0pt}{6pt}{3pt}
\titleformat{\subsection}{\normalsize\bfseries}{\thesubsection}{0.45em}{}
\titlespacing*{\subsection}{0pt}{5pt}{2pt}

\newfontfamily\zhfont{PingFang SC}
\newcommand{\blank}[1][8em]{\colorbox{yellow!35}{\rule{0pt}{1.1em}\hspace{#1}}}
\newcommand{\eg}{E_{\mathrm{g}}}
\newcommand{\ehull}{E_{\mathrm{hull}}}
\newcommand{\form}{x_{\mathrm{form}}}
\newcommand{\struc}{x_{\mathrm{struc}}}
\newcommand{\unc}{u}
\newcommand{\domain}{d_{\mathrm{dom}}}

\begin{document}

\begin{center}
{\Large \textbf{AI-Driven Design of Functional Boron Nitride Materials}}\\
{\normalsize Research Plan}
\end{center}

\textbf{Principal investigator, position and affiliated department}\\
Prof. Ma Chen, Assistant Professor, Department of Computer Science, City University of Hong Kong.

\textbf{Collaborator(s) and respective department(s), if any}\\
{\zhfont 深圳市中能天奕}.

\textbf{Duration of the project}\\
\blank[18em]

\section*{Abstract}
Boron nitride (BN) materials are promising for wide-band-gap electronics, dielectric layers, thermal management, ultraviolet optoelectronics, and two-dimensional heterostructures. This project aims to discover functional BN and BN-related materials using artificial intelligence. We will construct a BN-centered materials space, learn property models from public materials databases, optimize candidate formulas and structures under multi-objective targets, and prioritize candidates for structure generation and validation. The planned study will cover BN, BCN, BC$_2$N, SiBN, doped BN, and chemically related analogues. Expected outcomes include a curated BN data layer, AI property models, ranked candidate families, uncertainty-aware selection criteria, first-pass structure hypotheses, and a reusable research platform for academic and industrial collaboration.

\section*{Key objectives}
\begin{enumerate}
  \item \textbf{Define the BN design space.}
  \[
    \mathcal{C}_{\mathrm{BN}}=\{c:\{\mathrm{B},\mathrm{N}\}\subseteq
    \mathrm{elements}(c)\},\quad
    c\in\{\mathrm{BN},\mathrm{BCN},\mathrm{BC_2N},\mathrm{SiBN},\ldots\}.
  \]
  \item \textbf{Predict target properties.}
  \[
    f_{\theta}:(\form,\struc)\rightarrow
    \mathbf{y}=(\eg,\Delta H,\ehull,\unc,\rho_{\mathrm{app}}).
  \]
  \item \textbf{Rank and select candidates.}
  \[
    S(c)=w_1\hat{\eg}(c)-w_2\hat{\ehull}(c)-w_3\hat{\unc}(c)
    -w_4\domain(c)+w_5\rho_{\mathrm{BN}}(c)+w_6\rho_{\mathrm{novel}}(c).
  \]
  \item \textbf{Generate structure hypotheses.}
  \[
    c \mapsto \{s_{c,1},\ldots,s_{c,k}\}
    \mapsto \{\mathrm{geometry},\mathrm{relaxation},\mathrm{property\;recheck}\}.
  \]
\end{enumerate}

\section*{Research methodology}
\subsection*{1. Candidate universe}
The project starts from public 2D materials resources. 2DMatPedia provides an open database of more than 6,000 monolayer structures generated by top-down and bottom-up strategies \textsuperscript{[1]}. JARVIS-Tools supports programmatic access to atomistic materials datasets \textsuperscript{[2]}, while Matbench and matminer provide benchmark and featurization foundations for materials ML \textsuperscript{[3,4]}. Additional reference spaces include the Materials Project, C2DB, and related computational materials resources \textsuperscript{[5,6]}.

We will represent each record as
\[
  r_i=(c_i,s_i,\mathbf{y}_i,p_i),\qquad
  \mathcal{D}_{\mathrm{BN}}=\{r_i:c_i\in\mathcal{C}_{\mathrm{BN}}\},
\]
where $p_i$ records provenance. Candidate expansion will combine known BN motifs, same-group substitution, dopant enumeration, and collaborator-proposed chemical families.

\subsection*{2. AI models}
The modelling target is not a single score, but a design vector:
\[
  \max_{c,s}\; \mathbf{g}(c,s)=
  \big[\eg,\;-\ehull,\;\rho_{\mathrm{dielectric}},\;\rho_{\mathrm{UV}},
  \rho_{\mathrm{2D}},\;-\unc\big].
\]
Model families will include descriptor models, composition neural networks, crystal graph models, and structure-generation models, drawing on established materials-ML ideas such as CGCNN, ALIGNN, CrabNet, and CDVAE \textsuperscript{[7--10]}. Where suitable, pretrained interatomic-potential models such as CHGNet may support rapid structure screening \textsuperscript{[11]}.

\subsection*{3. Candidate selection}
The output for each candidate will be a decision vector
\[
  \mathbf{z}(c)=
  (\hat{\eg},\hat{\ehull},\hat{\unc},\domain,\rho_{\mathrm{BN}},
  \rho_{\mathrm{novel}},a),
\]
with action label
\[
  a\in\{\mathrm{control},\mathrm{priority},\mathrm{explore},\mathrm{hold}\}.
\]
This connects model prediction to research action: reference controls, high-priority validation, exploratory families, and candidates reserved for later evidence.

\subsection*{4. Iteration and validation}
The project will use grouped evaluation, BN-family evaluation, uncertainty analysis, and rank-stability metrics:
\[
  \mathrm{MAE},\ \mathrm{RMSE},\ R^2,\ \tau_{\mathrm{Kendall}},
  \rho_{\mathrm{Spearman}},\ P(c\in\mathrm{Top}\text{-}k).
\]
The current local project already provides a starting implementation for public-data ingestion, BN-themed candidate generation, model comparison, rank stability, and first-pass structure handoff. The proposed research will extend this base toward larger BN-relevant datasets, stronger AI models, and structure-level candidate design.

\section*{Significance and major deliverables}
BN is a high-value family for wide-band-gap and 2D materials research. h-BN is central to ultraviolet optoelectronics and layered-material platforms \textsuperscript{[12,13]}, and BN-related analogues create a broad design space for AI-guided discovery. The scientific question is:
\[
  \textit{Which BN-related compositions and structures should be investigated next?}
\]

\textbf{Major deliverables}
\begin{enumerate}
  \item BN materials dataset: formulas, structures, properties, provenance, and references.
  \item AI models for $\eg$, stability proxy, uncertainty, and application-readiness prediction.
  \item Candidate families: BN, BCN, BC$_2$N, SiBN, doped BN, and same-group analogues.
  \item Multi-objective ranking: $S(c)$, $\unc(c)$, $\domain(c)$, $\rho_{\mathrm{novel}}(c)$, and action labels.
  \item First-pass structure hypotheses and validation handoff files.
  \item Research outputs: candidate tables, figures, demo, technical report, and manuscript-ready results.
\end{enumerate}

\section*{Funding requested with budget breakdown}
\begin{center}
\small
\begin{tabular}{p{0.20\linewidth}p{0.51\linewidth}p{0.17\linewidth}}
\toprule
\textbf{Budget Item} & \textbf{Detail breakdown} & \textbf{Amount (HKD)} \\
\midrule
Staffing & Student helper / research assistant for data curation, experiments, documentation, and demo maintenance. & \blank \\
Equipment & Compute, storage, and structure-validation resources. & \blank \\
General expenses & Cloud credits, backup, software services, hosting, and data-management costs. & \blank \\
Conference & Presentation at AI-for-science, materials informatics, or applied computing venues. & \blank \\
Overseas Travel & Optional collaborator visit, conference travel, or partner technical meeting. & \blank \\
\midrule
 & \textbf{Total Funding} & \blank \\
\bottomrule
\end{tabular}
\end{center}

\section*{Ethics and safety}
The initial project is computational and data-driven. It uses public databases, software, models, and generated candidate structures. If later work involves chemical synthesis or regulated laboratory procedures, the responsible experimental party will handle the required safety approvals.

\section*{Selected references}
\small
\begin{enumerate}[leftmargin=1.2em,itemsep=1pt]
  \item Zhou, J. et al. 2DMatPedia: an open computational database of two-dimensional materials from top-down and bottom-up approaches. \textit{Scientific Data} 6, 86 (2019).
  \item Choudhary, K. et al. JARVIS-Tools: open-access software for atomistic data-driven materials design. \textit{npj Computational Materials} 6, 173 (2020).
  \item Dunn, A. et al. Benchmarking materials property prediction methods: the Matbench test set and Automatminer reference algorithm. \textit{npj Computational Materials} 6, 138 (2020).
  \item Ward, L. et al. Matminer: An open source toolkit for materials data mining. \textit{Computational Materials Science} 152, 60--69 (2018).
  \item Jain, A. et al. Commentary: The Materials Project: a materials genome approach to accelerating materials innovation. \textit{APL Materials} 1, 011002 (2013).
  \item Haastrup, S. et al. The Computational 2D Materials Database: high-throughput modeling and discovery of atomically thin crystals. \textit{2D Materials} 5, 042002 (2018).
  \item Xie, T. and Grossman, J. C. Crystal graph convolutional neural networks for an accurate and interpretable prediction of material properties. \textit{Physical Review Letters} 120, 145301 (2018).
  \item Choudhary, K. and DeCost, B. Atomistic line graph neural network for improved materials property predictions. \textit{npj Computational Materials} 7, 185 (2021).
  \item Wang, A. Y.-T. et al. CrabNet for explainable deep learning in materials science. \textit{npj Computational Materials} 7, 77 (2021).
  \item Xie, T. et al. Crystal diffusion variational autoencoder for periodic material generation. \textit{ICLR} (2022).
  \item Deng, B. et al. CHGNet as a pretrained universal neural network potential for charge-informed atomistic modelling. \textit{Nature Machine Intelligence} 5, 1031--1041 (2023).
  \item Watanabe, K., Taniguchi, T. and Kanda, H. Direct-bandgap properties and evidence for ultraviolet lasing of hexagonal boron nitride single crystal. \textit{Nature Materials} 3, 404--409 (2004).
  \item Cassabois, G., Valvin, P. and Gil, B. Hexagonal boron nitride is an indirect bandgap semiconductor. \textit{Nature Photonics} 10, 262--266 (2016).
\end{enumerate}

\end{document}
